
%!TEX root = ../hutbthesis_main.tex

\chapter{面向大规模集群的AirSim性能优化策略}

\section{优化总体目标与设计原则}

以解决20架以上无人机仿真性能断崖下跌为核心目标，具体各项目标如下：

1.帧率提升：30架仿真且FPS$\geq$30；

2.时延降低：平均时延$\leq$50ms；

3.内核占比降低：内核态CPU占比$\leq$30\%；

4.稳定性提升：连续运行1小时无崩溃、无失步；

5.可扩展性提升：支持30架以上平滑扩展。

（1）不破坏原有功能

优化方案严格保持对AirSim原有接口的兼容\textsuperscript{{[}5{]}{[}10{]}}，不改动现有API调用逻辑、ROS通信协议、传感器数据输出格式与物理仿真模型，确保原有项目、算法与调试流程可直接迁移使用。

（2）轻量化改造

以最小侵入方式修改底层源码，仅对线程调度、等待逻辑等核心瓶颈模块进行局部重构，不引入复杂第三方库或重型框架，保证代码简洁、编译稳定、维护成本低。

（3）可配置可切换

支持通过配置文件或运行时参数动态切换不同线程模型与同步策略，可在原生实现与多种优化方案之间灵活比对，便于实验测试、效果验证与场景适配。

（4）跨平台兼容

优化代码不依赖特定系统特性，可同时支持Windows与Linux操作系统，并兼容Unity与Unreal引擎环境，保证在不同平台下行为一致、性能稳定。

（5）开源可复用

整体设计遵循AirSim开源许可协议，代码结构规范、注释完整，成果可直接贡献至官方社区，便于其他研究者复用、扩展与二次开发。

本项目的优化方式如下图所示

\begin{figure}[H]
\centering
\includegraphics[width=0.88\linewidth,height=0.68\textheight,keepaspectratio]{images/image10.png}
\caption{优化方案总体设计架构图}
\label{fig:07}
\end{figure}

\section{线程等待模型优化设计}

针对AirSim低效等待机制的问题，设计了五种优化模型，如表3所示。

\begin{longtable}[]{@{}
  >{\raggedright\arraybackslash}p{(\columnwidth - 8\tabcolsep) * \real{0.1034}}
  >{\raggedright\arraybackslash}p{(\columnwidth - 8\tabcolsep) * \real{0.5674}}
  >{\raggedright\arraybackslash}p{(\columnwidth - 8\tabcolsep) * \real{0.1138}}
  >{\raggedright\arraybackslash}p{(\columnwidth - 8\tabcolsep) * \real{0.1063}}
  >{\raggedright\arraybackslash}p{(\columnwidth - 8\tabcolsep) * \real{0.1091}}@{}}
\caption{五种线程优化模型设计参数表}\\
\toprule\noalign{}
\begin{minipage}[b]{\linewidth}\raggedright
模型名称
\end{minipage} & \begin{minipage}[b]{\linewidth}\raggedright
核心实现
\end{minipage} & \begin{minipage}[b]{\linewidth}\raggedright
上下文切换开销
\end{minipage} & \begin{minipage}[b]{\linewidth}\raggedright
定时/响应精度
\end{minipage} & \begin{minipage}[b]{\linewidth}\raggedright
CPU占用率
\end{minipage} \\
\midrule\noalign{}
\endhead
\bottomrule\noalign{}
\endlastfoot
单线程分发器 &
采用集中式线程管理模式，通过统一调度器实现多无人机任务的集中管理与唤醒，避免多线程冗余竞争
& 极低 & 高 & 低 \\
标准阻塞睡眠 &
基于sleep\_for阻塞式等待机制，线程空闲时主动释放CPU资源，由系统内核统一调度唤醒
& 低 & 中 & 低 \\
协作式让出 &
通过yield函数实现线程主动让出CPU执行权，优先调度就绪状态的其他线程，保障线程间协作性
& 高 & 中 & 中 \\
改良自旋 &
结合自旋等待与阻塞睡眠的优势，采用sleep\_for(0)实现轻量化等待，平衡响应速度与资源占用
& 中 & 中高 & 中高 \\
纯自旋 &
采用无限空循环方式实现忙等，线程持续占用CPU资源不主动释放，无任何上下文切换操作
& 无 & 极高 & 极高 \\
\end{longtable}

注：(1)单线程分发器为本研究最优方案；(2)开销越低、精度越高，越适合大规模集群。

\begin{figure}[H]
\centering
\includegraphics[width=0.88\linewidth,height=0.68\textheight,keepaspectratio]{images/image11.png}
\caption{单线程分发器模型工作原理图}
\label{fig:08}
\end{figure}

核心思想：用一个专用线程管理所有无人机定时唤醒事件，避免大量线程并发等待。

优点：切换开销极低、内核占比低、精度高、扩展性好；

缺点：实现稍复杂。

\section{基于线程池的无人机管理与调度优化}

摒弃``一机一线程''模式，转而采用线程池：首先进行初始化，设定固定数量线程，该数量等同于CPU物理核心数，接着将无人机任务提交至线程池队列，通过线程复用，避免出现频繁创建与销毁线程的情况。线程池能够把线程数量从30以上降至8至16，大幅度减少竞争现象。其一为负载均衡，即将无人机任务均匀分散到线程池，防止单线程出现过载，其二是CPU亲和性，就是把线程绑定到指定的物理核心，减少核心之间的切换，提高缓存命中率。

依据集群规模动态调整等待间隔：若是小规模集群，采用高精度短间隔，若是大规模集群，则适度降低精度，以此保证整体的流畅度。

\section{\textbf{核心仿真模块的优化策略实施}}

AirSim核心修改文件如表4所示。

\begin{longtable}[]{@{}
  >{\raggedright\arraybackslash}p{(\columnwidth - 4\tabcolsep) * \real{0.1306}}
  >{\raggedright\arraybackslash}p{(\columnwidth - 4\tabcolsep) * \real{0.2865}}
  >{\raggedright\arraybackslash}p{(\columnwidth - 4\tabcolsep) * \real{0.5828}}@{}}
\caption{AirSim源码修改关键文件清单表}\\
\toprule\noalign{}
\begin{minipage}[b]{\linewidth}\raggedright
模块名称
\end{minipage} & \begin{minipage}[b]{\linewidth}\raggedright
文件名称
\end{minipage} & \begin{minipage}[b]{\linewidth}\raggedright
具体修改内容
\end{minipage} \\
\midrule\noalign{}
\endhead
\bottomrule\noalign{}
\endlastfoot
多旋翼管理 & MultiRotorConnector.cpp &
重构原生单机单线程架构，替换低效等待策略，集成优化后的线程调度与同步机制 \\
高精度定时器 & Timer.cpp &
改进底层定时逻辑，支持微秒级高精度等待，适配不同线程模型的时序同步需求 \\
仿真主循环 & SimModeBase.cpp &
优化集群仿真主循环流程，降低冗余调用开销，提升整体运行流畅度 \\
线程池管理 & ThreadPool.cpp &
新增轻量化线程池模块，实现线程复用、负载均衡与并发数限制，消除线程膨胀问题 \\
系统配置 & Settings.cpp &
新增优化策略配置项，支持运行时动态切换线程模型、等待方式与调度参数 \\
\end{longtable}

注：(1)所有修改为最小侵入式，不破坏原有功能；(2)改造模块均为性能瓶颈核心位置。

为了处理原生仿真主循环存在的冗余空转、定时精度不够、系统调用过于频繁等问题，着手对仿真主循环、定时器模块展开全面的重构工作。首先是对SimModeBase.cpp里的仿真主循环逻辑予以重构，把循环体内部那些无效的空转等待、冗余的条件判断、重复的函数调用进行梳理后移除，以此精简主循环的执行流程，降低不必要的CPU资源消耗，进而提高主循环的执行效率。其次是优化Timer.cpp中的定时器实现，借助底层硬件时钟机制达成微秒级的高精度定时，替换掉原生的低效定时方式，减少定时过程里的系统调用次数，降低内核态的开销，保障多机协同仿真的时序一致性。第三点是支持定时周期动态配置，能够借助配置文件或者运行时参数，对定时器周期进行灵活调整，以此适配不同的集群规模、不同仿真精度的需求，进而增强优化方案的场景适配能力。

为了提高优化方案的灵活性、可测试性、兼容性，设计并实现了灵活的编译配置、运行时切换机制。在编译阶段，借助宏定义参数设置，可以自由选择启用不同的线程优化模型，方便编译不同版本的仿真程序，以适配不同的测试需求，在运行阶段，支持通过JSON配置文件动态调整优化参数，达成不同等待模型、调度策略的运行时切换，无需重新编译源码，从而提高实验测试效率。严格保持AirSim原生默认运行模式，没有对其原有功能实施具有破坏性的修改，能够在同一个仿真环境里达成原生模式与优化模式之间的快速切换，如此一来方便对比测试两种模式的性能差异，进而为验证优化效果创造便捷有利的条件。

\section{优化方案的兼容性与可行性验证}

优化方案严格依照AirSim原生架构设计原则来执行，只是针对线程调度、定时器等核心模块进行局部优化工作，并没有对无人机物理模型、传感器数据输出、场景渲染等核心业务逻辑做出改动，以此保证与原有的API接口、ROS通信协议实现完全兼容。经由实验测试能够发现，优化后的系统能够正常加载原有的仿真场景，还可以解析传感器数据，并且能够与第三方算法工具实现无缝对接，充分满足科研、工程应用的需求。为了能够精确验证优化效果，采用控制变量法，在相同的硬件环境、相同的仿真场景之下，对比优化前后的核心性能指标，具体的数据如下表表5所示：

\begin{longtable}[]{@{}
  >{\raggedright\arraybackslash}p{(\columnwidth - 6\tabcolsep) * \real{0.2087}}
  >{\raggedright\arraybackslash}p{(\columnwidth - 6\tabcolsep) * \real{0.2087}}
  >{\raggedright\arraybackslash}p{(\columnwidth - 6\tabcolsep) * \real{0.3533}}
  >{\raggedright\arraybackslash}p{(\columnwidth - 6\tabcolsep) * \real{0.2293}}@{}}
\caption{性能指标测试结果表}\\
\toprule\noalign{}
\begin{minipage}[b]{\linewidth}\raggedright
性能指标
\end{minipage} & \begin{minipage}[b]{\linewidth}\raggedright
优化前（原生）
\end{minipage} & \begin{minipage}[b]{\linewidth}\raggedright
优化后（线程池+阻塞等待）
\end{minipage} & \begin{minipage}[b]{\linewidth}\raggedright
提升幅度
\end{minipage} \\
\midrule\noalign{}
\endhead
\bottomrule\noalign{}
\endlastfoot
平均响应延迟 & 87ms & 12ms & 86.2\% \\
上下文切换次数 & 120次/分钟 & 18次/分钟 & 85.0\% \\
内存占用率 & 42\% & 18\% & 57.1\% \\
仿真帧率 & 18fps & 42fps & 133.3\% \\
\end{longtable}

注：(1)指标覆盖实时性、资源开销、精度与稳定性；(2)合格标准依据无人机实时仿真行业要求。

由上表数据可知，优化后系统在响应速度、资源利用效率等方面均实现显著提升：平均响应延迟降低86.2\%，上下文切换次数减少85.0\%，内存占用率下降57.1\%，仿真帧率提升133.3\%，充分验证了线程优化、定时器重构等方案的有效性，解决了原有架构下的性能瓶颈，实现了``高效、稳定、可扩展''的设计目标。

\section{本章小结}

本章明确优化目标与原则，设计五种线程等待模型，实现线程池化、负载均衡、亲和性调度、定时器重构等优化策略，完成AirSim核心源码改造。优化方案具备高性能、高兼容、高稳定、可扩展等优点，为后续实验测试奠定坚实基础。
